\documentclass[a4paper,11pt]{article}

\usepackage[utf8]{inputenc}
\usepackage[T1]{fontenc}
\usepackage[ngerman]{babel}
\usepackage{amsmath,amssymb}
\usepackage{booktabs}
\usepackage{graphicx}
\usepackage{geometry}
\usepackage{hyperref}
\usepackage{tikz}
\usepackage{caption}
\usepackage{enumitem}
\usepackage{pifont}
\usetikzlibrary{arrows.meta, positioning, fit, backgrounds}

\geometry{margin=2.5cm}
\hypersetup{colorlinks=true, linkcolor=blue!50!black, citecolor=black, urlcolor=blue!50!black}

\title{\textbf{Digitaler Zwilling für die thermische Simulation \\
eines elektrodynamischen Levitators (TEAM~28-Benchmark)} \\[0.3em]
\large Konzept, Methodik und erwartete Ergebnisse}
\author{Fachgebiet TEMF \\ Projektbericht: Thermal Digital Twin}
\date{Stand: August 2026}

\begin{document}
\maketitle

\begin{abstract}
\noindent
Dieser Bericht beschreibt Idee, physikalische Methodik und Softwarearchitektur
eines \emph{digitalen Zwillings} (Digital Twin) für die thermische Vorhersage
am elektrodynamischen Levitator nach dem TEAM-28-Benchmark (COMPUMAG). Statt
kommerzieller CAE-Software wurde ein selbst entwickelter Python-Löser gewählt,
weil das eigentliche Ziel keine einmalige Simulation, sondern ein
\emph{echtzeitfähiges} Vorhersagemodell ist, das auf einem Smartphone per
AR-Ansicht läuft. Der Bericht erläutert die physikalische Modellkette
(Elektromagnetik~$\to$~Joulsche Verlustwärme~$\to$~Wärmeleitung), die
numerische Umsetzung (FEM-Offline-Lösung plus echtzeitfähiges Reduced-Order-Model),
die Entwicklungsmethodik (KI-gestützte, agentische Softwareentwicklung mit
strenger physikalischer Verifikation) sowie die bisherigen und erwarteten
Ergebnisse. Offene Fragen und Grenzen des aktuellen Modells werden transparent
benannt.
\end{abstract}

\tableofcontents
\clearpage

% ---------------------------------------------------------------------------
\section{Einleitung und Motivation}
% ---------------------------------------------------------------------------

\subsection{Die Aufgabenstellung}
Das TEAM-Problem~28 ist ein etablierter Benchmark der
Computational-Electromagnetics-Gemeinschaft (COMPUMAG) zur Validierung
selbst geschriebener elektromagnetischer Löser: eine leitfähige Aluminiumscheibe
levitiert über einer Spulenanordnung, die mit Wechselstrom (hier $\hat{\imath}
\approx 5\,\mathrm{A}$, $f = 50\,\mathrm{Hz}$) betrieben wird. Der zeitlich
veränderliche magnetische Fluss induziert Wirbelströme in der Scheibe, die
sowohl eine Hubkraft (Levitation) als auch eine Verlustwärme (Erwärmung)
erzeugen. Der am Fachgebiet TEMF aufgebaute Versuchsstand entspricht diesem
Benchmark strukturell, weist aber wesentliche Zusatzelemente auf, die im
Originalproblem nicht vorkommen: einen ferromagnetischen Kern
($\mu_r \approx 1000$) und einen ferromagnetischen Außenring, beide mit
messbarer Eigenerwärmung.

\subsection{Vom Simulationsmodell zum digitalen Zwilling}
Die Aufgabe lautet nicht, \emph{einmal} ein Temperaturfeld zu berechnen,
sondern ein Modell zu bauen, das
\begin{itemize}[nosep]
  \item das Temperaturfeld $T(r,z,t)$ \textbf{kontinuierlich} vorhersagt, während
        Nutzerinnen und Nutzer den Betriebsstrom in Echtzeit ändern,
  \item gegen den realen Versuchsstand \textbf{validiert} wird (IR-Kamera,
        Thermoelemente),
  \item am Ende über einen \textbf{QR-Code} auf einem Smartphone in
        Augmented Reality (AR) betrachtet werden kann.
\end{itemize}
Diese drei Anforderungen zusammen definieren einen \emph{digitalen Zwilling}
im eigentlichen Sinn -- nicht nur ein Simulationsergebnis, sondern ein
laufendes, validiertes, interaktives Modell des physischen Systems. Genau
diese Anforderung ist der Ausgangspunkt für alle in diesem Bericht
begründeten methodischen Entscheidungen.

% ---------------------------------------------------------------------------
\section{Physikalisches Grundkonzept: Warum Echtzeit möglich ist}
% ---------------------------------------------------------------------------

Ein digitaler Zwilling, der bei jeder Stromänderung eine volle FEM-Simulation
neu rechnet, wäre auf einem Smartphone nicht realisierbar (Rechenzeit:
Sekunden bis Minuten pro Lauf). Die Lösung liegt in einer physikalischen
Beobachtung, nicht in einer Software-Abkürzung:

\begin{quote}
\itshape
Bei fester Frequenz ($50\,\mathrm{Hz}$) und linearen Materialeigenschaften
skalieren alle Joulschen Verluste quadratisch mit dem Strom
($P \propto \hat{\imath}^2$), während sich die \textbf{räumliche
Verteilung} der Verluste \emph{nicht} ändert -- nur ihre Amplitude.
\end{quote}

Daraus folgt eine zweistufige Architektur:
\begin{enumerate}[nosep]
  \item \textbf{Offline (einmalig):} Die elektromagnetische FEM wird
        \emph{einmal} bei einem Referenzstrom $\hat{\imath}_{\mathrm{ref}}$
        gelöst. Ergebnis ist eine ortsaufgelöste Verlustkarte
        $\hat{q}(r,z)$ (Wirbelstromverluste in Scheibe und Eisen) sowie
        Spulenverluste $\hat{P}_{\mathrm{Cu}}$.
  \item \textbf{Online (Echtzeit, Millisekunden):} Für einen beliebigen
        Betriebsstrom $\hat{\imath}$ und die zuletzt bekannte
        Körpertemperatur $\bar T$ wird die Verlustkarte nur noch
        \textbf{skalar multipliziert}:
        \[
          q(r,z;\hat{\imath},\bar T) = \hat q(r,z)\cdot
          \left(\frac{\hat{\imath}}{\hat{\imath}_{\mathrm{ref}}}\right)^{2}
          \cdot \frac{\sigma(\bar T)}{\sigma(T_{\mathrm{ref}})}.
        \]
        Diese Verlustleistung speist ein reduziertes thermisches Modell
        (Reduced-Order-Model, ROM), das ebenfalls nur mit einfacher
        Vektor-/Skalar-Arithmetik arbeitet.
\end{enumerate}

Diese Zerlegung ist der eigentliche methodische Kern des Projekts: Die
"`Intelligenz“" steckt vollständig im einmaligen Offline-Löser; die
Echtzeitfähigkeit entsteht rein aus der physikalischen Struktur des Problems
(Linearität, feste Frequenz), nicht aus einer vereinfachten oder
"`gecheateten“" Physik.

% ---------------------------------------------------------------------------
\section{Physikalisches Modell}
% ---------------------------------------------------------------------------

Das Modell besteht aus einer Kette von drei gekoppelten Teilproblemen, die
alle achsensymmetrisch sind und daher in der reduzierten 2D-Ebene $(r,z)$
gelöst und für die Darstellung zu 3D rotiert werden.

\subsection{Elektromagnetik (Phasor-Formulierung)}
Da nur die \emph{zyklusgemittelte} Verlustwärme für die (im Vergleich zu
$50\,\mathrm{Hz}$ sehr langsame) Thermik relevant ist, wird das
elektromagnetische Feld nicht zeitschrittweise bei $50\,\mathrm{Hz}$
simuliert, sondern als komplexer Phasor $A_\varphi(r,z)$ (azimutale
Komponente des Vektorpotentials) gelöst:
\[
  -\nabla\!\cdot\!\left(\nu \nabla A_\varphi\right)
  + \nu\,\frac{A_\varphi}{r^{2}}
  + j\omega\sigma\, A_\varphi = J_s ,
  \qquad \nu = \frac{1}{\mu_r \mu_0},\quad \omega = 2\pi f .
\]
Hierbei ist $J_s$ die eingeprägte Stromdichte in den Spulen (Innenspule und
Außenspule mit entgegengesetztem Wicklungssinn), $\sigma$ die elektrische
Leitfähigkeit (Aluminium, Kupfer, Eisen) und $\nu$ die magnetische
Reluktivität, die im ferromagnetischen Kern/Ring stark abfällt und so den
magnetischen Fluss bündelt.

\subsection{Joulsche Verlustwärme (zyklusgemittelt)}
Aus dem Phasorfeld folgt die induzierte Wirbelstromdichte
$J_e = -j\omega\sigma A_\varphi$ und daraus die zeitlich gemittelte
Verlustdichte
\[
  q(r,z) = \frac{|J_e|^{2}}{2\sigma}
         = \tfrac{1}{2}\,\sigma\,\omega^{2}\,|A_\varphi|^{2}
         \quad [\mathrm{W/m^3}].
\]
Der Faktor $\tfrac12$ folgt aus der Mittelung von $\sin^2(\omega t)$ über
eine Periode. Diese quadratische Abhängigkeit von $|A_\varphi|$ ist die
physikalische Grundlage der $I^2$-Skalierung aus Abschnitt~3.

\subsection{Temperaturabhängige Leitfähigkeit}
Da eine Erwärmung um $70\,\mathrm{K}$ den spezifischen Widerstand um rund
$27\,\%$ ändert, wird dieser Effekt nicht vernachlässigt:
\[
  \sigma(T) = \frac{\sigma_0}{1+\alpha (T-T_0)}, \qquad
  \alpha \approx 3.9\times 10^{-3}\,\mathrm{K^{-1}}.
\]
Bemerkenswert ist die \textbf{gegenläufige} Wirkung: Die Scheibenverluste
sind proportional zu $\sigma_{\mathrm{Al}}$ (heißere Scheibe $\to$ geringere
Verluste, negative Rückkopplung), während die Spulenverluste proportional
zu $1/\sigma_{\mathrm{Cu}}$ sind (heißere Spule $\to$ höhere Verluste,
positive Rückkopplung). Da beide Effekte nur die \emph{Amplitude} und nicht
das \emph{räumliche Muster} ändern, bleibt die Zweistufenarchitektur aus
Abschnitt~3 gültig -- $\sigma(T)$ wird als weiterer Laufzeit-Skalar
eingerechnet, ohne die FEM neu zu lösen.

\subsection{Wärmeleitung}
Die Verlustdichte $q$ dient als Quellterm der instationären
Fourier-Gleichung:
\[
  \rho c_p \frac{\partial T}{\partial t} = \nabla\!\cdot\!(k\nabla T) + q,
\]
mit Robin-Randbedingung für Konvektion,
$-k\,\partial T/\partial n = h\,(T-T_\infty)$. Im stationären Fall muss die
Energiebilanz exakt erfüllt sein:
\[
  \int_V q\,\mathrm dV \;=\; \oint_{\Gamma} h\,(T-T_\infty)\,\mathrm dA .
\]
Diese Bilanzgleichung ist der zentrale Verifikationstest des thermischen
Lösers (siehe Abschnitt~7).

% ---------------------------------------------------------------------------
\section{Numerische Umsetzung und Softwarearchitektur}
% ---------------------------------------------------------------------------

\subsection{Diskretisierung}
Beide Teilprobleme (Elektromagnetik, Wärmeleitung) werden mit der
\textbf{Finite-Elemente-Methode (FEM)} in Galerkin-Formulierung diskretisiert,
mit linearen Dreieckselementen (P1) auf einem strukturierten 2D-Netz
$(r,z)$. Die achsensymmetrische Schwachform trägt das Volumengewicht
$2\pi r$:
\[
  \int k\,\nabla T\!\cdot\!\nabla v \cdot 2\pi r \,\mathrm dA
  + \oint_\Gamma h\,T\,v\cdot 2\pi r\,\mathrm ds
  = \int p\,v\cdot 2\pi r\,\mathrm dA
  + \oint_\Gamma h\,T_\infty\,v\cdot 2\pi r\,\mathrm ds .
\]
Für das elektromagnetische Teilproblem ergibt die Assemblierung eine
\textbf{komplexe} dünnbesetzte Matrix (wegen des $j\omega\sigma$-Terms), die
direkt mit \texttt{scipy.sparse.linalg.spsolve} gelöst wird. Bei
$\sim\!10^{4}$ Unbekannten (2D, achsensymmetrisch) genügt ein direkter
Löser -- ein iteratives Verfahren oder eine 3D-Vernetzung wäre für diese
Problemgröße unnötiger Mehraufwand.

Ein wichtiges Detail zur Netzauflösung: Die Skintiefe in Aluminium bei
$50\,\mathrm{Hz}$ beträgt $\delta \approx 12\,\mathrm{mm}$, deutlich mehr als
die Scheibendicke von $3\,\mathrm{mm}$. Der Wirbelstrom ist somit über die
Dicke näherungsweise gleichverteilt -- eine feine Vernetzung in
$z$-Richtung ist nicht erforderlich; die Netzverfeinerung konzentriert sich
stattdessen auf den Spalt zwischen Spule und Scheibe sowie auf die
Außenkanten.

\subsection{Reduced-Order-Model (ROM) für Echtzeit}
Aus der einmaligen FEM-Lösung wird ein niederdimensionales, echtzeitfähiges
Modell abgeleitet:
\begin{itemize}[nosep]
  \item \textbf{$I^2$-Skalierung + Zeitkonstante erster Ordnung}
        (\texttt{rom.py}) für die Scheibe, mit thermischer Zeitkonstante
        $\tau \approx 4{,}07\,\mathrm{min}$ bei $R=80\,\mathrm{mm}$.
  \item \textbf{Verteiltes RC-Netzwerk} (\texttt{twin\_core.py}) für Spulen,
        Eisenkern/-ring und den gemeinsamen Luftknoten -- jeder Körper als
        Wärmekapazität $C=\rho c_p V$, verbunden über Leitwerte $G$ bzw.
        Konvektionsleitwerte $hA$. Dieses Netzwerk ist bewusst so gewählt,
        weil es sich direkt an realen Sensordaten kalibrieren lässt
        (Abschnitt~7).
\end{itemize}
Der gesamte Online-Schritt reduziert sich auf einfache Vektor-/
Skalararithmetik ($O(n)$), lauffähig in Millisekunden im Browser eines
Smartphones -- exakt die Voraussetzung für die AR-Anwendung.

\subsection{Softwarepipeline}
Die Implementierung folgt einer strikten Ein-Weg-Pipeline, wobei jede Datei
genau eine Verantwortung trägt und alle physikalischen Parameter zentral in
\texttt{params.yaml} liegen (\emph{single source of truth}, keine
Konstante ist im Code fest verdrahtet):

\begin{center}
\begin{tikzpicture}[
  node distance=6mm and 10mm,
  box/.style={draw, rounded corners, align=center, minimum width=3.1cm,
              minimum height=1.05cm, font=\footnotesize},
  arr/.style={-Stealth, thick}
]
\node[box] (params) {\texttt{params.yaml}\\[-2pt]\tiny (SSOT aller Parameter)};
\node[box, below=of params] (config) {\texttt{config.py}\\[-2pt]\tiny Einheiten, $I_{\mathrm{peak}}$};
\node[box, below left=8mm and -3mm of config] (em) {\texttt{em\_solver.py}\\[-2pt]\tiny EM-FEM, Verlustkarte};
\node[box, below right=8mm and -3mm of config] (th) {\texttt{thermal\_solver.py}\\[-2pt]\tiny Wärme-FEM, Bilanz};
\node[box, below=16mm of config] (rom) {\texttt{rom.py} / \texttt{twin\_core.py}\\[-2pt]\tiny Echtzeit-ROM (SSOT-Integrator)};
\node[box, below=of rom] (out) {\texttt{build\_twin\_html\_fem.py}\\[-2pt]\tiny Bake $\to$ AR-HTML + QR};
\node[box, right=14mm of rom] (sensor) {\texttt{data\_io.py} + Arduino\\[-2pt]\tiny Sensor-Kalibrierung};

\draw[arr] (params) -- (config);
\draw[arr] (config) -- (em);
\draw[arr] (config) -- (th);
\draw[arr] (em) -- (rom);
\draw[arr] (th) -- (rom);
\draw[arr] (rom) -- (out);
\draw[arr] (sensor) -- (rom);
\end{tikzpicture}
\captionof{figure}{Vereinfachte Softwarepipeline: von der Parameterdatei über
die Offline-FEM-Löser zum echtzeitfähigen ROM und der finalen
AR-Auslieferung. Der Sensorpfad speist die Kalibrierung des ROM.}
\end{center}

Eine \emph{Kreuzvalidierung} (\texttt{xval\_twin.py}, Playwright-basiert)
stellt sicher, dass der in JavaScript "`gebackene“" Echtzeit-Integrator der
AR-HTML-Datei exakt mit dem Python-Referenzintegrator (\texttt{twin\_core.py})
übereinstimmt -- eine Fehlerklasse, die sonst leicht unbemerkt entsteht
(Koeffizient im Python-Code korrigiert, aber die HTML-Datei nicht neu
gebacken).

% ---------------------------------------------------------------------------
\section{Warum kein kommerzielles CAE-Werkzeug}
% ---------------------------------------------------------------------------

Werkzeuge wie SimScale, COMSOL/ANSYS Maxwell oder FEMM könnten das
elektromagnetisch-thermische Teilproblem grundsätzlich lösen. Sie erfüllen
jedoch nicht die eigentliche Anforderung -- ein echtzeitfähiges,
sensor-kalibrierbares, auf einem Smartphone lauffähiges Modell:

\begin{table}[h]
\centering\small
\begin{tabular}{@{}lccccc@{}}
\toprule
Kriterium & SimScale & COMSOL/Maxwell & FEMM & Elmer/FEniCS & \textbf{Python (gewählt)} \\
\midrule
Niederfrequenz-EM (Wirbelstrom) & begrenzt & \checkmark & \checkmark & \checkmark & \checkmark (validiert an TEAM~28) \\
EM$\to$Thermik-Kopplung & schwierig & \checkmark & schwach & komplex & \checkmark direkt in-memory \\
ROM / Echtzeitfähigkeit & \ding{55} & teilweise, teuer & \ding{55} & \ding{55} & \checkmark \textbf{Kernentwurf} \\
Export nach Web/AR & \ding{55} & \ding{55} & \ding{55} & \ding{55} & \checkmark 1 HTML-Datei \\
Sensor-Kalibrierung & \ding{55} & proprietär & \ding{55} & selbst zu bauen & \checkmark \texttt{data\_io.py} \\
Läuft auf macOS & \checkmark & \checkmark (teuer) & \ding{55} nur Windows & \checkmark & \checkmark \\
Kosten & Core-Hour-Limit & tausende~€ & kostenlos & kostenlos & kostenlos \\
Black Box? & ja & ja & teilweise & nein & \textbf{nein -- volle Kontrolle} \\
\bottomrule
\end{tabular}
\caption{Vergleich der Werkzeugoptionen gegen die Anforderungen des digitalen
Zwillings (nicht gegen eine reine Einzelsimulation).}
\end{table}

Der entscheidende Punkt ist strukturell, nicht nur praktisch: Kommerzielle
und viele Open-Source-FEM-Pakete kapseln den Löser als Black Box -- jede
Stromänderung erfordert einen vollständigen Solver-Neustart. Um die in
Abschnitt~3 beschriebene $I^2$-Struktur auszunutzen (FEM einmal lösen,
danach nur noch skalieren), muss der Löser auf Quellcode-Ebene
kontrollierbar sein: Zwischenergebnisse (die Verlustkarte $\hat q(r,z)$
\emph{vor} der Multiplikation mit $I^2$) müssen gezielt herausgezogen und
weiterverwendet werden können. FEMM entfällt zusätzlich aus einem rein
praktischen Grund (nur Windows-nativ, Entwicklungsrechner ist macOS).

% ---------------------------------------------------------------------------
\section{Entwicklungsmethodik: KI-gestützte, agentische Softwareentwicklung}
% ---------------------------------------------------------------------------

\subsection{Ansatz}
Für die Implementierung wurde bewusst ein \textbf{KI-gestützter, agentischer
Entwicklungsprozess} gewählt: Ein KI-Coding-Agent (Claude Code) arbeitet als
Pair-Programmer, der iterativ Code schreibt, physikalische
Verifikationsläufe ausführt und die Ergebnisse interpretiert, während die
inhaltliche und physikalische Verantwortung -- Modellannahmen,
Parameterwerte, Validierungskriterien, Freigabe jeder Änderung -- vollständig
beim studentischen Team verbleibt. Diese Arbeitsweise unterscheidet sich von
klassischem "`Copy-Paste aus einem Chat-Fenster“": Der Agent hat Zugriff auf
das gesamte Repository, führt Code tatsächlich aus, prüft Testkriterien und
dokumentiert jede Änderung nachvollziehbar in der Versionsgeschichte (Git).

\subsection{Warum dieser Ansatz für dieses Projekt passt}
Das Projekt hat einen stark \textbf{explorativen} Charakter: Geometrie und
Materialparameter des realen Aufbaus (z.\,B. die Permeabilität des
Eisenkerns) waren zu Projektbeginn nicht bekannt und mussten schrittweise
durch Messung, Nachrechnung und Kalibrierung eingegrenzt werden. Das führt
zu vielen kurzen Iterationszyklen: Parameter ändern
$\to$ FEM neu lösen $\to$ mit Messdaten vergleichen $\to$ Modell anpassen.
Ein KI-gestützter Agent verkürzt genau diese Zyklen erheblich, weil er
Routineaufgaben (Refactoring, Testläufe, Konsistenzprüfungen zwischen
mehreren Dateien, Aktualisieren von Dokumentation) parallel zur inhaltlichen
Arbeit der Studierenden übernimmt.

\subsection{Wie die physikalische Korrektheit trotzdem sichergestellt wird}
Ein KI-Agent kann -- wie jede Softwarequelle -- Fehler einführen. Die
Qualitätssicherung basiert deshalb \textbf{nicht auf Vertrauen in den
Agenten}, sondern auf harten, automatisch prüfbaren physikalischen
Kriterien, die für jede Änderung erneut ausgeführt werden müssen:
\begin{itemize}[nosep]
  \item \textbf{Energiebilanz}: $\int q\,\mathrm dV = \oint h(T-T_\infty)\,\mathrm dA$
        muss auf $0{,}000\,\%$ Abweichung stimmen.
  \item \textbf{$I^2$-Konsistenz}: $P(2\hat\imath)/P(\hat\imath)$ muss exakt
        $4{,}000000$ ergeben.
  \item \textbf{Kreuzvalidierung} (\texttt{xval\_twin.py}): der in
        JavaScript gebackene Echtzeit-Integrator muss bitgenau mit dem
        Python-Referenzintegrator übereinstimmen.
  \item \textbf{Kalibrierung gegen reale Messdaten}: IR-Kamera- und
        Thermoelementdaten (Abschnitt~7) als externe, vom Code unabhängige
        Wahrheitsquelle.
  \item \textbf{Vollständige Versionsgeschichte}: jede Änderung, jede
        Parameteranpassung und jede Fehlerkorrektur ist in Git und im
        Projekt-Changelog nachvollziehbar dokumentiert -- inklusive
        rückgängig gemachter oder als falsch erkannter Ansätze.
\end{itemize}
Dieses Vorgehen entspricht dem etablierten Prinzip \emph{"`Verifikation vor
Fertigstellungsbehauptung“"}: Eine Änderung gilt erst dann als korrekt, wenn
das zugehörige physikalische Kriterium nachweislich erfüllt ist -- nicht,
weil der Code plausibel aussieht oder eine KI ihn produziert hat.

\subsection{Vorteile der Gesamtmethodik}
\begin{itemize}[nosep]
  \item \textbf{Geschwindigkeit bei gleichbleibender Nachvollziehbarkeit}:
        Iterationszyklen (Messung $\to$ Modellanpassung $\to$ erneute
        Validierung) verkürzen sich deutlich, ohne die physikalische
        Kontrolle zu verlieren, da jede Iteration weiterhin an den harten
        Kriterien aus 6.3 gemessen wird.
  \item \textbf{Konsistenz über viele gekoppelte Dateien}: Physik wird an
        zwei Stellen implementiert (Python-Referenz und gebackenes
        JavaScript für die AR-Ansicht); Änderungen an beiden Stellen
        synchron zu halten ist eine Routineaufgabe, die sich gut delegieren
        lässt, solange die Kreuzvalidierung als Prüfinstanz bestehen bleibt.
  \item \textbf{Dokumentationsdisziplin}: Jede Kalibrierungsrunde,
        jede Parameteränderung und jede offene Frage wird laufend
        mitgeschrieben (\texttt{CHANGELOG.md}) -- eine Nebenwirkung des
        agentischen Arbeitsablaufs, die für die spätere Nachvollziehbarkeit
        des Projekts wertvoll ist.
  \item \textbf{Kein Zielkonflikt mit dem Kursziel}: Da alle physikalischen
        Annahmen, Gleichungen und Kalibrierungen von den Studierenden
        definiert, geprüft und verantwortet werden, bleibt der
        Lernanspruch des Kurses -- Verständnis von der schwachen Form bis
        zur Matrix -- vollständig erhalten; die KI übernimmt Schreibarbeit,
        nicht physikalisches Verständnis.
\end{itemize}

% ---------------------------------------------------------------------------
\section{Validierung und bisherige Ergebnisse}
% ---------------------------------------------------------------------------

\subsection{Interne Konsistenzprüfungen}
\begin{itemize}[nosep]
  \item Energiebilanz des thermischen Lösers: Abweichung $0{,}000\,\%$.
  \item $I^2$-Skalierungstest: $P(2\hat\imath)/P(\hat\imath) = 4{,}000000$
        (exakt).
  \item Gebietsgrößenvalidierung (nach Vorgabe des Fachgebiets): Vergleich
        Dirichlet- vs.\ Neumann-Randbedingung auf einem $1\times1\,\mathrm m$
        Rechengebiet -- alle Abweichungen $<1\,\%$ (größte: Scheibenverlust
        $0{,}767\,\%$), Gebiet somit ausreichend groß gewählt.
  \item Sättigungsprüfung im Eisen: $B_{\max}\approx 0{,}66\,\mathrm T$,
        deutlich unterhalb der Sättigungsgrenze ($\approx 1{,}5\,\mathrm T$)
        -- die Annahme eines linearen $\mu_r$ ist im Betriebspunkt gültig.
\end{itemize}

\subsection{Aktuelles quantitatives Ergebnis (î\,=\,5\,A, $T_\infty=20\,^\circ$C)}
\begin{table}[h]
\centering
\begin{tabular}{@{}lr@{}}
\toprule
Verlustquelle & $P$\,[W] \\
\midrule
Wirbelstromverlust Scheibe (Al) & 25{,}81 \\
Wirbelstromverlust Eisenkern + -ring & 5{,}86 \\
Ohmscher Spulenverlust (innen 52{,}32 + außen 54{,}12) & 106{,}44 \\
\midrule
\textbf{Gesamt} & \textbf{138{,}11} \\
\bottomrule
\end{tabular}
\caption{Verlustleistungen bei Nennbetrieb, reproduziert nach der
Neuvermessung der Geometrie vom Juli 2026.}
\end{table}

\subsection{Kalibrierung gegen reale Messdaten}
Zwei reale Betriebspunkte wurden per IR-Kamera (HIKMICRO) vermessen
($190\,\mathrm V\to5\,\mathrm A$ als Hauptbetriebspunkt,
$270\,\mathrm V\to7{,}8\,\mathrm A$ als Kalibrieranker). Das
RC-Netzwerkmodell der Spulen wurde durch Rückwärtsrechnung
(\texttt{refit\_hA.py}) direkt gegen den vollständigen, nichtlinearen
Zeitintegrator kalibriert und trifft den stationären Zustand exakt
($T_{\mathrm{innen}}=79{,}00\,^\circ$C, $T_{\mathrm{außen}}=74{,}00\,^\circ$C
bei $7{,}8\,\mathrm A$). Die Scheiben- und Kerntemperaturen aus der
IR-Aufnahme gelten aufgrund unbekannter Emissivität von blankem Aluminium
als nicht belastbar -- nur die dunkel lackierten Spulen liefern verlässliche
IR-Werte; die verbleibende Validierung von Scheibe und Kern ist Aufgabe der
geplanten Kontaktsensorik (Abschnitt~7.4).

\subsection{Sensorpfad für zukünftige Kalibrierung}
Der Software-seitige Kalibrierpfad
(\texttt{arduino/thermal\_sensor.ino}~$\to$~\texttt{data\_io.py}~$\to$
\texttt{rom.calibrate\_UA()}) ist vollständig implementiert und mit einer
synthetischen Messreihe (\texttt{mock\_sensor\_data.csv}) durchgehend
getestet. Es fehlt ausschließlich der reale Sensoraufbau (Arduino +
2$\times$ MAX31855-Thermoelement-Wandler + IR-Referenzmessung), dessen
Einkaufsliste dem Fachgebiet zur Beschaffung vorliegt.

% ---------------------------------------------------------------------------
\section{Erwartete Ergebnisse und Ausblick}
% ---------------------------------------------------------------------------

\subsection{Kurzfristig erwartete Ergebnisse}
\begin{itemize}[nosep]
  \item Ein vollständig kalibrierter, echtzeitfähiger digitaler Zwilling,
        dessen Spulentemperaturen mit realer IR-Messung übereinstimmen
        (bereits erreicht, RMS-Fehler $\approx 2{,}5$--$3\,^\circ$C).
  \item Eine reale Aufheiz- \emph{und} Abkühlkurve von Spule und Scheibe
        (bislang liegt nur eine Aufheizkurve vor), womit das bislang nicht
        eindeutig bestimmbare RC-Netzwerk (insbesondere der gemeinsame
        Luftknoten) vollständig identifizierbar wird.
  \item Eine erste belastbare Kontaktmessung der Scheibenunterseite, die
        klärt, ob im stationären Zustand tatsächlich die Spule oder die
        Scheibe die höhere Temperatur erreicht (im Modell aktuell nur
        unentschieden abgebildet).
  \item Eine veröffentlichte AR-Ansicht (\texttt{digital\_twin\_fem.html})
        mit zugehörigem QR-Code, aufrufbar ohne Installation auf einem
        Smartphone.
\end{itemize}

\subsection{Mittelfristig}
\begin{itemize}[nosep]
  \item Eine dichtere Kalibriertabelle Stelldrehwinkel~$\to$~Strom (aktuell
        nur drei Ankerpunkte).
  \item Eine reale Strommessvorrichtung, sodass der digitale Zwilling einen
        tatsächlich gemessenen Zeitverlauf $I(t)$ statt eines konstanten
        Werts verarbeitet.
  \item Eine belastbarere Bestimmung der Eisenpermeabilität (aktuell ein
        Platzhalterwert $\mu_r=1000$ ohne gemessene B-H-Kurve), um die unten
        genannte offene Frage zur Levitationshöhe zu klären.
\end{itemize}

% ---------------------------------------------------------------------------
\section{Offene Fragen und Grenzen des aktuellen Modells}
% ---------------------------------------------------------------------------

Im Sinne wissenschaftlicher Transparenz werden die derzeit ungelösten Punkte
explizit benannt, statt sie zu verschweigen:

\begin{itemize}[nosep]
  \item \textbf{Levitationshöhe (Hubkraft) stimmt nicht mit der Beobachtung
        überein.} Das Modell sagt eine Hubkraft voraus, die deutlich größer
        ist als die Gewichtskraft, mit einer rechnerischen
        Gleichgewichtshöhe von $\approx 11{,}7$--$14{,}7\,\mathrm{mm}$,
        während der reale, sichtbare Spalt nur $7$--$8\,\mathrm{mm}$
        beträgt. Magnetische Sättigung wurde als Ursache
        ausgeschlossen (Kraftänderung $<0{,}1\,\%$ bei nichtlinearer
        Rechnung); wahrscheinlichste Ursache ist der unvermessene
        $\mu_r$-Platzhalterwert des Eisens. Eine willkürliche
        Rückwärtsanpassung von $\mu_r$ an die beobachtete Höhe wurde
        bewusst \emph{nicht} vorgenommen, um die Messdaten nicht zu
        verfälschen.
  \item \textbf{Eisenparameter unvermessen}: $\mu_r=1000$ und
        $\sigma\approx1{,}0\times10^6\,\mathrm{S/m}$ sind plausible
        Platzhalter für "`weichstahlähnliches“" Material, keine gemessenen
        Werte.
  \item \textbf{Original-TEAM-28-Benchmark (ohne Eisen, 20\,A)} liefert
        aktuell $z_{\mathrm{eq}}\approx14{,}5\,\mathrm{mm}$ gegenüber
        $11{,}3\,\mathrm{mm}$ laut Referenztabelle ($\approx28\,\%$
        Abweichung) -- pausiert, da nicht blockierend für den realen
        Versuchsstand, aber ebenfalls offen.
  \item \textbf{RC-Netzwerk teilweise unteridentifiziert}: Ohne eine reale
        Abkühlkurve lässt sich der gemeinsame Luftknoten nicht eindeutig
        von den übrigen Parametern trennen (siehe Abschnitt~8.1).
\end{itemize}

% ---------------------------------------------------------------------------
\section{Zusammenfassung}
% ---------------------------------------------------------------------------

Das Projekt verfolgt einen physikalisch begründeten Weg zu einem
echtzeitfähigen digitalen Zwilling: eine einmalige, hochauflösende FEM-Lösung
der Elektromagnetik- und Wärmeleitungsgleichungen wird über die
$I^2$-Skalierung und ein kalibrierbares Reduced-Order-Model auf
Millisekunden-Laufzeit reduziert. Diese Struktur ist mit gekapselten
kommerziellen CAE-Werkzeugen nicht realisierbar und rechtfertigt die
Entscheidung für eine selbst kontrollierte Python-Implementierung. Die
Entwicklung wurde durch einen KI-gestützten, agentischen Arbeitsablauf
beschleunigt, dessen Ergebnisse jedoch ausschließlich über harte,
automatisierte physikalische Kriterien (Energiebilanz, $I^2$-Konsistenz,
Kreuzvalidierung, Abgleich mit realen Sensordaten) freigegeben werden --
nicht über Vertrauen in die KI selbst. Erste Kalibrierungsergebnisse an der
Spulentemperatur bestätigen die Tragfähigkeit des Modells; die
Levitationshöhe bleibt als offene physikalische Frage bestehen und wird mit
der geplanten Sensorhardware weiter untersucht.

\end{document}
